\paragraph*{04.05.06.01 Qualitätsmanagementplan für das Projekt entwickeln, die Implementierung
überwachen und gegebenenfalls überarbeiten}
\kcilabel{para:04.05.06.01_Qualitätsmanagementplan}{04.05.06.01}
\label{para:04.05.06.01_Qualitätsmanagementplan}

% Task
\par Wie im \ref{para:04.03.02.04_Supportfunktionen} geschildert war die Sicherstellung der Qualität der \gls{ricefw}-Entwicklung von zentraler Bedeutung, um die erfolgreiche Umsetzung der manigfaltigen regulatorischen Anforderungen, die Einhaltung von SAP-Standards und die Erfüllung der Stakeholder-Erwartungen in Form von qualitativ hochwertigen Ergebnissen zu gewährleisten. 

% Action
\par Zur Sicherstellung der Qualität entwickelte ich zusammen mit dem Programmtestmanager einen umfassenden Qualitätsmanagementplan, der die spezifischen Qualitätsanforderungen des \gls{wsr} adressierte.\endnote{Im \gls{s4} agiere ich als stv. \gls{tm}} Dieser Plan umfasste die Definition von Qualitätsstandards, die Festlegung von Qualitätskontrollprozessen und die Implementierung von Prozessen zur kontinuierlichen Verbesserung (vgl. \autoref{fig:TestpläneSST} und \autoref{fig:UATDatenDoku}).

% Result
\par Wie im \ref{para:04.04.05.02_OwnershipUndCommitment} dargestellt, wurde der Qualitätsmanagementplan erfolgreich umgesetzt und hat zu einer exzellenten Qualität der \gls{ricefw}-Entwicklung im \gls{epa} und \gls{zp1} in Welle 1 geführt.

% Reflect
\par Die Umsetzung des Qualitätsmanagementplans hat zu einer signifikanten Verbesserung der Qualität der \gls{ricefw}-Entwicklung geführt und die Stakeholder-Zufriedenheit erhöht.